% Katalog I: Inventar — jede API-Einheit einmal, mit Codestelle (Slot 9).
% Der visuelle Teil (A–H) zeigt, WIE etwas aussieht. Diese Sektion zeigt,
% WAS es ist und WO es lebt — die zwei Fragen, die beim Aussortieren
% unmittelbar auf die Entscheidung folgen.
\ZSFNewColumn
\StartChapter[kat:inventar]{I · Inventar}

\begin{runintext}
  Jede Einheit der stabilen API genau einmal, nach Quelldatei gruppiert.
  Ein Strich statt einer ID heisst: hat keine eigene Darstellung, steht
  deshalb nicht im visuellen Teil.
\end{runintext}

\SubsectionBar[kat:inv-arten]{Drei Arten von Eingriff}
\ZSFindex{Inventar}

\begin{ZSFlist}
  \ZSFItem{Umgebung}{Eigener Name mit \ZSFlabel{begin/end}. Streichen heisst:
  Name weg, vorhandene Aufrufe auf einen anderen Baustein umstellen.}
  \ZSFItem{Reglerwert}{Ein Zweig in der Box-Fabrik. Streichen heisst: der
  Wert entfällt, die Vorbelegung bleibt — kein Aufruf verliert seine Box.}
  \ZSFItem{Makro}{Ein Befehl. Streichen heisst: Befehl weg, Aufrufe ersetzen.}
\end{ZSFlist}
\Zweck{Die Unterscheidung entscheidet, was ein »weg damit« überhaupt kostet.}

\begin{defbox}[Wo das meiste liegt][weight=quiet]
  \ZSFhl{34 der 104 Namen stehen in 60\_boxes.tex} — mit 78 KB zugleich die
  grösste Datei. Wer dort etwas entfernt, arbeitet an der Box-Fabrik, die
  alle Boxen gemeinsam tragen.
\end{defbox}

\SubsectionBar[kat:inv-boxes]{60\_boxes.tex · 34}
\ZSFindex{Box-Fabrik}

\begin{tablebox}[Umgebungen]
  \begin{ZSFtable}[font=dense, rows=tight, colsep=tight]{Y{1.4} Z{0.6}}
    \ZSFheaderRow \ZSFhead{Name} & \ZSFhead{ID} \\
    defbox & B-02 \\
    figbox & B-08 \\
    formulabox & B-06 \\
    goalbox & B-09 \\
    runintext & B-01 \\
    splitbox & B-05 \\
    tablebox & B-07 \\
    valuegrid & D-14 \\
    warnbox & B-03 \\
    ZSFboxgroup & B-10 \\
    ZSFinset & B-11 \\
    ZSFlist & B-04
  \end{ZSFtable}
\end{tablebox}

\begin{tablebox}[Makros aus derselben Datei]
  \begin{ZSFtable}[font=dense, rows=tight, colsep=tight]{Y{1.4} Z{0.6}}
    \ZSFheaderRow \ZSFhead{Name} & \ZSFhead{ID} \\
    GoalCondition & G-09 \\
    GoalStep & G-09 \\
    GoalTarget & G-09 \\
    SubsectionBar & A-04 \\
    SubsubsectionBar & A-06 \\
    formulanote & G-07 \\
    ZSFBoxesBreakableOn & --- \\
    ZSFBoxesBreakableOff & --- \\
    ZSFBreakReservesOff & --- \\
    ZSFDerivationCase & G-10 \\
    ZSFItem & E-14 \\
    ZSFNewColumn & A-08 \\
    ZSFTitleHeader & A-01 \\
    ZSFTitleTag & E-10 \\
    ZSFdanger & E-05 \\
    ZSFfig & F-01 \\
    ZSFfigcaption & F-05 \\
    ZSFfigside & F-03 \\
    ZSFformulaline & G-06 \\
    ZSFgroupcols & B-10 \\
    ZSFimage & D-09 \\
    ZSFsep & E-12
  \end{ZSFtable}
\end{tablebox}

\SubsectionBar[kat:inv-knobs]{Reglerwerte · 18 Schlüssel}
\ZSFindex{Reglerwert}

\begin{tablebox}[Alle in der Box-Fabrik]
  \begin{ZSFtable}[font=dense, rows=tight, colsep=tight]{Y{1.1} Y{1.3} Z{0.6}}
    \ZSFheaderRow \ZSFhead{Schlüssel} & \ZSFhead{ausser Default}
                  & \ZSFhead{ID} \\
    align & center & C-12 \\
    atomic & Flag & C-15 \\
    bodyparskip & inherit & C-14 \\
    colsep & tight & D-07 \\
    font & dense & C-13 \\
    frame & strong, hard, none & C-09 \\
    grid & horizontal, none & D-10 \\
    header & false & D-02 \\
    ordered & Flag & C-20 \\
    padx & bar, none & C-05 \\
    pady & tight, none & C-07 \\
    part & first, mid, last & H-05 \\
    rows & roomy, tight & D-05 \\
    split & Anteil & C-16 \\
    splitalign & center, bottom & C-17 \\
    tone & neutral, warn & C-01 \\
    weight & quiet & C-03 \\
    zebra & false & D-03
  \end{ZSFtable}
\end{tablebox}
\Zweck{Vier gehören genau einer Box: ordered, grid, part und die vier
Tabellen-Regler. Alle übrigen gelten auf jeder Box.}

\SubsectionBar[kat:inv-rest]{Die übrigen Module}
\ZSFindex{Modul}

\begin{tablebox}[10\_math.tex · 19]
  \begin{ZSFtable}[font=dense, rows=tight, colsep=tight]{Y{1.4} Z{0.6}}
    \ZSFheaderRow \ZSFhead{Name} & \ZSFhead{ID} \\
    R, C, N, Z, Q & G-01 \\
    dd, vect, abs, norm & G-02 \\
    ZSFsumAuto, ZSFlimAuto & G-03 \\
    ZSFbraceunder, -over & G-04 \\
    ZSFmhlA bis D & G-08 \\
    SetZSFsumMode & --- \\
    SetZSFlimMode & ---
  \end{ZSFtable}
\end{tablebox}

\begin{tablebox}[20\_tables.tex · 12]
  \begin{ZSFtable}[font=dense, rows=tight, colsep=tight]{Y{1.4} Z{0.6}}
    \ZSFheaderRow \ZSFhead{Name} & \ZSFhead{ID} \\
    ZSFtable & D-01 \\
    ZSFheaderRow, ZSFhead & D-01 \\
    ZSFspan & D-08 \\
    L, C, R, F & D-12 \\
    Y, Z, Q & D-13 \\
    SetZSFzebraBG & ---
  \end{ZSFtable}
\end{tablebox}
\Zweck{C, Q, R und Z heissen hier Spaltentyp und in 10\_math Zahlmenge —
derselbe Buchstabe, zwei Bedeutungen, getrennt durch den Kontext.}

\begin{tablebox}[50\_typography\_semantics.tex · 13]
  \begin{ZSFtable}[font=dense, rows=tight, colsep=tight]{Y{1.4} Z{0.6}}
    \ZSFheaderRow \ZSFhead{Name} & \ZSFhead{ID} \\
    ZSFkeyword & E-01 \\
    ZSFlabel & E-03 \\
    ZSFhl & E-04 \\
    ZSFconclusion & E-06 \\
    ZSFref & E-07 \\
    ZSFsectionref & E-08 \\
    ZSFScriptRef & E-09 \\
    ZSFfontDiagramLabel & E-17 \\
    ZSFfontDiagramLabelSmall & E-17 \\
    ZSFkeywordStyle & --- \\
    ZSFlabelStyle & --- \\
    ZSFhlStyle & --- \\
    ZSFDiagramLabelScale & ---
  \end{ZSFtable}
\end{tablebox}

\begin{tablebox}[30\_layout\_spacing.tex · 15]
  \begin{ZSFtable}[font=dense, rows=tight, colsep=tight]{Y{1.4} Z{0.6}}
    \ZSFheaderRow \ZSFhead{Name} & \ZSFhead{ID} \\
    textVorBox, textNachBox & E-15 \\
    ZSFgap & --- \\
    ZSFSectionGap & --- \\
    ZSFBodyFontSize & --- \\
    ZSFDensityFactor & --- \\
    ZSFDensityBoxes & --- \\
    ZSFDensityText & --- \\
    ZSFDensityTables & --- \\
    ZSFDensityStructure & --- \\
    ZSFLeadingFactor & --- \\
    ZSFBarGapFactor & --- \\
    ZSFBreakReserveFactor & --- \\
    ZSFChapterDensity & --- \\
    ZSFChapterDensityReset & ---
  \end{ZSFtable}
\end{tablebox}

\begin{tablebox}[40\_colors\_structure.tex · 6]
  \begin{ZSFtable}[font=dense, rows=tight, colsep=tight]{Y{1.4} Z{0.6}}
    \ZSFheaderRow \ZSFhead{Name} & \ZSFhead{ID} \\
    StartChapter & A-02 \\
    StartFrontChapter & A-03 \\
    ZSFDeclareQuantity & E-16 \\
    ZSFDeclareTone & --- \\
    ZSFQuantityColorsOn & --- \\
    ZSFQuantityColorsOff & ---
  \end{ZSFtable}
\end{tablebox}

\begin{tablebox}[55\_readability.tex · 5]
  \begin{ZSFtable}[font=dense, rows=tight, colsep=tight]{Y{1.4} Z{0.6}}
    \ZSFheaderRow \ZSFhead{Name} & \ZSFhead{ID} \\
    ZSFbreak & --- \\
    ZSFnobreak & --- \\
    ZSFallowbreak & --- \\
    ZSFReadableBodyOn & --- \\
    ZSFReadableBodyOff & ---
  \end{ZSFtable}
\end{tablebox}

\begin{tablebox}[Opt-in-Module]
  \begin{ZSFtable}[font=dense, rows=tight, colsep=tight]{Y{1.1} Y{1.3} Z{0.6}}
    \ZSFheaderRow \ZSFhead{Modul} & \ZSFhead{Name} & \ZSFhead{ID} \\
    11 & Operatoren & H-01 \\
    11 & drawbrace, tikzmark & H-02 \\
    12 & pgfplots & H-03 \\
    65 & codebox & H-04 \\
    67 & CodeLine & H-06 \\
    67 & InlineComment & H-07 \\
    67 & OverlineComment & H-08 \\
    66 & ZSFindex & H-09 \\
    66 & ZSFindexsee & H-10 \\
    66 & ZSFIndexShowPageNumber & ---
  \end{ZSFtable}
\end{tablebox}
